\section{Citation system, evidence policy and exclusions}\label{sec:rules}

\subsection{Two kinds of reference}\label{sec:citations}
\begin{enumerate}
\item \textbf{Theory} uses numbered bibliography entries in text, for example ``the discrete Karush--Kuhn--Tucker conditions are the baseline of sparse active-set optimisation~\cite{wohlmuth2011contact,hintermuller2002pdas}''. The bibliography is in journal style. The file \code{plan/refs.bib} starts from \code{paper/refs.bib} and adds the contact and active-set references.
\item \textbf{Code} uses a code anchor that names the module and the symbol and the release: \anchor{fea_engine.contact_solvers.solve_contact_static_sparse @ v1.0.1}. The anchor uses the importable name, not a file path, so it survives file moves; the back matter maps every anchor to its path and to the repository URL. The package versions at the time of writing are \fea{} 1.0.1 and \rom{} 0.2.0.
\end{enumerate}

\subsection{Evidence policy}
\begin{itemize}
\item Every number in the text is generated by a script, a test or a result file of the repository, and the caption names it. As in the paper, numbers are macros produced by a table-generation script, so that text and results cannot drift apart.
\item Timings are labelled with the machine and are single runs unless stated; no test asserts a speed.
\item A claim that was not measured is written as not measured. Limits are placed next to the result they qualify.
\item Items marked \texttt{to produce} in \Cref{sec:chapters} are created before the chapter is written, and a chapter is not finished while it contains one.
\end{itemize}

\subsection{What is not written}
\begin{itemize}
\item No commercial solver workflows, file formats or meshes; every step is native to the from-scratch implementation.
\item No user-interface code, web servers or extensive plot customisation; figures come from short scripts and the text concerns the numerical arrays.
\item No elementary content: basic matrix transformations, first shape-function derivations and 1-D truss examples are skipped, except where a patch test needs one.
\item No wrapped high-level optimisation or machine-learning packages; the book relies on NumPy, SciPy and PyTorch.
\end{itemize}

\subsection{Inventory of items still to produce}
\begin{table}[ht]
\centering\small
\begin{tabularx}{\linewidth}{@{}l X l@{}}
\toprule
Chapter & Item & Kind \\
\midrule
1 & Architecture diagram; dispatch-overhead script and table & figure, script, table \\
2 & \code{meshio} round-trip check; grading figure & script, figure \\
3 & Locking comparison of \code{Hex8Solid3D}, \code{Hex8SolidBbar}, \code{Hex20Solid3D} & figure \\
4 & Newmark order-of-accuracy table; explicit stability-limit test & script, table \\
6 & Speedup figure from \code{batched_kernels.json}; chunking memory table & figure, table \\
8 & Offline and online cost table from experiment e1 & table \\
10 & Certified-bound figure & figure \\
11 & Interface size against accuracy table from experiment e3 & table \\
14 & Summary table of the five reproductions & table \\
\bottomrule
\end{tabularx}
\caption{Items that do not exist yet and are needed by the chapters.}\label{tab:todo}
\end{table}
